% !TEX root =../tfmETSI.tex
\chapter{Introduction}\label{chp:introduction}
\epigraph{Learning directly from raw text about images is a promising alternative which leverages a much broader source of supervision.}{Radford et al., 2021}

\section{Motivation}
\label{sec:intro-motivation}

Histopathology --- the microscopic examination of stained tissue sections --- is the gold-standard test for the definitive diagnosis of cancer and many other diseases~\cite{vanderlaak2021deeplearning}. Automating any meaningful part of the pathologist's workflow (triaging slides, flagging regions of interest, suggesting a candidate diagnosis) is one of the highest-impact targets for medical AI, and the visual patterns involved are well-defined enough to be machine-learnable in principle.

The practical bottleneck is data. Labelled histopathology is expensive to produce --- it takes pathologist time and clinical context that do not scale --- and patient privacy further restricts what can be shared. This makes any approach that relies on broad, label-cheap supervision particularly attractive.

Foundation models~\cite{bommasani2021foundation} are designed for exactly this regime: models trained on broad data at large scale that can then be adapted to a wide range of downstream tasks with little task-specific supervision. In the vision--language line, CLIP~\cite{radford2021clip} is the canonical example and the paradigm under study in this work. In plain terms, CLIP trains two encoders side by side, one for images and one for text, so that matched image--text pairs land close together in a shared embedding space. Once trained, the model can be applied to a new classification task simply by writing a short text prompt for each class, with no further training required. We unpack the architecture and training signal in Chapter~\ref{chp:methodology}; for now it is enough to know that CLIP is the closest thing to a general-purpose vision--language model available off the shelf.

Within the same paradigm, the histopathology-specific variant we use in this work is \textbf{PLIP}~\cite{huang2023plip}: the same dual-encoder contrastive recipe, trained on $\sim 208$k pathology image--text pairs scraped from medical Twitter. PLIP is one of several histopathology foundation models available (the broader family is reviewed in Chapter~\ref{chp:stateofart}), and we pick it here because it is the closest single-variable swap from our baseline: same CLIP architecture, ViT-B compute footprint, and openly released weights. PLIP plays two roles in this work: as a zero-shot reference for what an off-the-shelf domain-pretrained model can do without any task-specific training, and as the histopathology-pretrained image encoder swapped into the main pipeline in Chapter~\ref{chp:results}, where it produces the headline finding.

\textbf{The question this project asks is where the pretraining of a CLIP-style histopathology pipeline actually pays off, and where it does not} --- evaluated across head-side, loss-side, text-side, and image-side interventions, both in-distribution and on three external datasets covering two organs. The framing is deliberate: CLIP as the object of study rather than the classifier itself, so the findings should generalise to other scientific domains where short or missing labels are the norm.


\section{Project Scope Evolution}
\label{sec:intro-scope}

The project began with a different scope. The initial direction was semantic segmentation of histopathology slides --- pixel-level identification of tissue regions --- which is closer to what a clinical workflow actually produces. That direction was abandoned within the first few weeks: annotated segmentation data at the scale we needed was too scarce, and producing it would have required pathologist time well beyond the project budget. We pivoted to image-level classification on the same domain, which preserves the methodological questions about CLIP and histopathology while removing the data-cost wall.

The training source is LC25000~\cite{borkowski2019lc25000}, a 5-class lung-plus-colon patch dataset that fits the frozen-backbone training regime without needing pixel-level labels. External evaluation runs on three datasets chosen to span shift severity and two organs: NCT-CRC-HE-7K~\cite{kather2019nct} (close-shift colon), Chaoyang~\cite{zhu2021chaoyang} (different-hospital colon), and LungHist700~\cite{tabatabaei2024lunghist} (different-source lung). The study is organised as a series of single-variable comparisons: each experiment changes exactly one component of a fixed baseline (frozen ResNet50 + frozen BERT with a CLIP head trained on top) --- the image encoder, the text encoder, the input preprocessing, or the classifier head --- and holds the rest identical. The delta on each external dataset then traces back to one specific design decision, and the comparisons that follow in Chapter~\ref{chp:results} tell us which of those decisions actually matter.


\section{Research Questions}
\label{sec:intro-rq}

The work is structured around six research questions. \textbf{The central one is cross-cutting}: how does a CLIP-style histopathology pipeline generalise to data the model has never seen? The other five each isolate a single axis of the pipeline (baseline behaviour, the CLIP-vs-non-CLIP architecture choice, input preprocessing, the text encoder, and the image encoder) so that the contribution of each can be measured separately when answering this cross-cutting question.

\begin{itemize}
    \item \textbf{RQ1} --- What baseline classification and generalisation behaviour does a frozen general-purpose CLIP pipeline (ImageNet ResNet50 + BERT) achieve on LC25000, and what failure modes appear?
    \item \textbf{RQ2} --- Does the CLIP paradigm validate its \emph{foundational-model} premise relative to a non-CLIP classifier on the same image encoder, both in-distribution and on data the model was never trained on?
    \item \textbf{RQ3} --- Does stain normalisation help when we evaluate on datasets the model was never trained on?
    \item \textbf{RQ4} --- Does using a text encoder pretrained on biomedical text (BioBERT~\cite{lee2020biobert}, PubMedBERT~\cite{gu2021pubmedbert}), instead of generic BERT~\cite{devlin2019bert}, improve performance? And does the prompt format (bare class name vs longer composed prompt) interact with the encoder choice?
    \item \textbf{RQ5} --- Does the image encoder's \emph{pretraining choice} matter? Specifically: replacing the general-purpose ResNet50~\cite{he2016resnet} with PLIP's histopathology-pretrained ViT-B/32~\cite{huang2023plip,dosovitskiy2021vit}, and with DINOv2's general-domain self-supervised ViT-B/14~\cite{oquab2024dinov2} (a strong off-the-shelf image encoder trained without labels on $\sim 142$M general natural images) as a control that isolates ``ViT architecture and self-supervision'' from ``histopathology pretraining specifically''.
    \item \textbf{RQ6} --- How well does each variant generalise to data the model has never seen, evaluated across three external datasets that span shift severity and two organs?
\end{itemize}


\section{Contributions}
\label{sec:intro-contributions}

The central contribution of this project is a controlled measurement of how much each major design choice in a CLIP-style histopathology pipeline actually matters, with every claim framed as a delta over the same baseline --- frozen ResNet50 + frozen BERT, single fixed prompt, stratified 80/10/10 split on LC25000 --- and iterated through one experiment at a time in Chapter~\ref{chp:results}.

Headline findings (full results in Chapter~\ref{chp:results}):
\begin{itemize}
    \item \textbf{Image encoder pretraining is the dominant lever.} Replacing ResNet50 with PLIP's pretrained ViT-B/32 (histopathology image--text contrastive) gains $+0.067$ F1 on close-shift colon (NCT-CRC), $+0.027$ on different-hospital colon (Chaoyang), and $+0.335$ on different-source lung (LungHist700). No other intervention tested moves the needle comparably.
    \item \textbf{The lung-side advantage decomposes into three additive steps of comparable magnitude.} The DINOv2 and PLIP zero-shot controls together decompose PLIP's $+0.335$ F1 lung advantage into three contributions: ViT + self-supervised pretraining on general images ($+0.152$), histopathology pretraining specifically ($+0.087$), and LC25000 supervision on top ($+0.096$).
    \item \textbf{Five other interventions move the needle far less, or only in narrow regimes.} Stain Macenko (NCT-NORM reference) is the only input-side lever that produces a measurable positive gain under shift ($+0.030$ on Chaoyang, $+0.020$ on LungHist700), but its magnitude is dominated by the image-encoder lever on the same datasets and it regresses on NCT-CRC ($-0.022$, which arrives Macenko-normalised at release). Biomedical text encoders (BioBERT, PubMedBERT) tie or trail baseline across all 6 grid cells on NCT, and the best biomedical cell (PubMedBERT + name\_only) extended to Chaoyang and LungHist700 confirms no transfer benefit on any of the three datasets. A non-CLIP plain classifier ties CLIP in-distribution and on close-shift external data but loses on heavier shift; prompt-format variations within the text-encoder comparison are within noise.
    \item \textbf{CLIP's foundational-model premise is validated under shift, not in-distribution.} CLIP and a plain CE classifier tie at saturation in-distribution and on close-shift external; CLIP's edge appears under heavier shift via the bounded-cosine + temperature-scaled classifier geometry, which does not degenerate the way high-confidence softmax does.
\end{itemize}

Methodological contributions:
\begin{itemize}
    \item Single-variable-change protocol: each experiment differs from the baseline in exactly one design choice, so the resulting delta is attributable.
    \item Per-organ external reporting: results are presented per dataset, not averaged; averaging across datasets would hide the cross-organ regime change observed in this study.
    \item Reproducibility: artifact bundle (fixed-seed split, cached DINOv2 features, weights, metric files, plots) versioned alongside the report in the project repository.
    \item Explicit written discussion of dataset and methodological limits (Chapter~\ref{chp:conclusion}).
\end{itemize}


\section{Document Structure}
\label{sec:intro-structure}
\begin{itemize}
    \item Chapter~\ref{chp:introduction} (Introduction) --- motivation, scope, research questions, contributions.
    \item Chapter~\ref{chp:stateofart} (State of the Art) --- foundation models, multimodal contrastive lineage, pathology FMs, where this work fits.
    \item Chapter~\ref{chp:methodology} (Methodology) --- CLIP architecture, image and text encoders, datasets, preprocessing, baseline, evaluation protocol, reproducibility.
    \item Chapter~\ref{chp:results} (Results and Discussion) --- the experiments and what each measured, ending in a synthesis section.
    \item Chapter~\ref{chp:conclusion} (Conclusions and Future Work) --- one-paragraph-per-RQ conclusions, limitations, future work.
\end{itemize}
